% !TEX program = lualatex
% Auto-generated from YAML (v3) — 02: TeamsとFormsの基本操作

\documentclass[aspectratio=43,professionalfonts,handout,10pt]{beamer}
\usepackage{beamer_template}

\newcommand{\LectureNum}{02}
\newcommand{\LectureTitle}{TeamsとFormsの基本操作}
\newcommand{\CourseName}{情報リテラシー}
\newcommand{\TextPages}{-}
\newcommand{\NextTopic}{情報とメディアの特性}
\newcommand{\NextPages}{pp.2-5}

\begin{document}
% --- Slide 1: title ---

\begin{frame}{\CourseName\quad 第\LectureNum 回「\LectureTitle 」}
  {\small \TermName\quad \TeacherName\quad ／\quad 教科書 \TextPages}

  \vfill

  \begin{block}{今日の流れ}
    \begin{enumerate}
      \item 個人Microsoftアカウント（outlook.jp）の作成
      \item TeamsとFormsの使い分け
      \item Microsoft Teamsの基本操作
      \item デジタルコミュニケーションのマナー
    \end{enumerate}
  \end{block}
  \vfill

  \begin{block}{今日の目標}
    \begin{itemize}
      \item Microsoft365などのWebシステムを理解し，取り扱うことができ，情報発信に利用できる。
    \end{itemize}
  \end{block}
  \vfill

  \begin{exampleblock}{キーワード}
    \small \textbf{Teams} ／ \textbf{チャネル} ／ \textbf{Forms} 
  \end{exampleblock}
\end{frame}

% --- Slide 2: practice (演習1統合) ---

\begin{frame}{演習1：個人Microsoftアカウント（outlook.jp）の作成}
  {\small PCを卒業後も使えるよう，今のうちに個人アカウントを作成しておく。}

  \vfill

  \begin{block}{手順（各自スマートフォンで実施）}
    \begin{enumerate}
      \item スマートフォンのブラウザで https://signup.live.com にアクセスする
      \item 「新しいメールアドレスを取得」をタップし，好きなアドレス（例：yamada.taro.2025@outlook.jp）を入力する
      \item パスワードを設定する（16文字以上，英字＋数字を含む）
      \item 姓名・生年月日を入力する
      \item SMSまたはメールで届く確認コードを入力して本人確認を完了する
      \item アカウント作成完了後，メールアドレスとパスワードをノートにメモする
    \end{enumerate}
  \end{block}
\end{frame}

% --- Slide 3: Teams vs Forms 比較表 ---

\begin{frame}{TeamsとFormsの使い分け}
  \begin{block}{ツールの特徴}
    \begin{tabular}{lll}
      \hline
      & \textbf{Teams} & \textbf{Forms} \\
      \hline
      役割 & コミュニケーション & 回答の収集・集計 \\
      主な機能 & チャット・通話・ファイル共有 & アンケート・テスト作成 \\
      用途 & やりとり・連絡・議論 & 出欠・感想・確認テスト \\
      \hline
    \end{tabular}
  \end{block}

  \vfill

  \begin{exampleblock}{使い分けのポイント}
    \begin{itemize}
      \item やりとりや連絡 → \textbf{Teams}
      \item 回答の収集や集計 → \textbf{Forms}
    \end{itemize}
  \end{exampleblock}
\end{frame}

% --- Slide 4: exercise ---

\begin{frame}{演習2：隣の人にチャットを送ってみよう}
  \begin{block}{手順}
    \begin{enumerate}
      \item 隣の席の人と名前を教え合う
      \item Teamsの左側メニューから「チャット」をクリックする
      \item 画面上部の「新しいチャット」（ペンのアイコン）をクリックする
      \item 「宛先」に隣の人の名前を入力し，候補から選択する（見つからない場合は学籍番号で検索する）
      \item メッセージ欄に「こんにちは！○○です。よろしくお願いします。」と入力し，送信する
      \item 相手からの返信を確認する
    \end{enumerate}
  \end{block}

  \vfill

  \begin{exampleblock}{ポイント}
    チャット＝1対1やグループの会話，チャネル投稿＝クラス全体への発信
  \end{exampleblock}
\end{frame}

% --- Slide 5: exercise ---

\begin{frame}{演習3：会議に参加する}
  {\small Teamsの会議機能を使って，オンライン会議に参加する方法を学ぶ。}

  \vfill

  \begin{block}{手順}
    \begin{enumerate}
      \item Teamsの情報リテラシーの会議から「参加」ボタンをクリックする
      \item カメラ・マイクのオフを確認し，「今すぐ参加」をクリックする
      \item チャット欄にメッセージを送信する
      \item 「退出」ボタンで会議を終了する
    \end{enumerate}
  \end{block}

\end{frame}

% --- Slide 6: マナー ---

\begin{frame}{デジタルコミュニケーションのマナー：アリ？ナシ？}
  \begin{block}{次の行動，あなたはどう思う？}
    \begin{enumerate}
      \item チャネルに毎日「おはよう！」と投稿する
      \item 間違って投稿した内容をそのまま放置する
      \item 「@全員」を何度も使って連絡する
      \item リアクションだけで返信し，言葉を添えない
      \item 件名を「質問」だけにしてメールを送る
    \end{enumerate}
  \end{block}

  \vfill

  \begin{exampleblock}{マナーの基本}
    \begin{itemize}
      \item 投稿は簡潔・丁寧に，スレッド返信を活用する
      \item \texttt{@全員} は本当に必要なときだけ
      \item 迷ったら一言添えると誤解が減る
    \end{itemize}
  \end{exampleblock}
\end{frame}

% --- Slide 7: exercise ---

\begin{frame}{演習4（ケーススタディ）}
  

  \vfill

  \begin{block}{以下の場面では Teams・Forms どちらが適切か}
    \begin{enumerate}
      \item クラス全員から今日の授業の感想を集めたい
      \item 緊急連絡をすぐに全員に伝えたい
      \item 出席確認を自動で集計したい
    \end{enumerate}
  \end{block}
\end{frame}

% --- Slide 8: チェックリスト ---

\begin{frame}{まとめ・チェックリスト}
  \begin{block}{自己チェック（できたら $\checkmark$ をつけよう）}
    \begin{itemize}
      \item[$\square$] outlook.jpの個人Microsoftアカウントを作成し，メモした
      \item[$\square$] TeamsとFormsの使い分けを説明できる
      \item[$\square$] Teamsでチャットを送り，会議に参加した
      \item[$\square$] デジタルコミュニケーションのマナーを理解した
    \end{itemize}
  \end{block}

  \vfill

  {\small 質問があれば授業後またはTeamsで受け付けます。}
\end{frame}

\end{document}
